\documentclass[12pt]{article}
\usepackage[utf8]{inputenc}
\usepackage[margin=1in]{geometry}   % 设置页边距
\usepackage{amsmath, amssymb}       % 数学符号与公式
\usepackage{enumitem}               % 自定义列表
\usepackage{unicode-math}
\usepackage[english]{babel}         % 设置整个文档为英文
\usepackage{graphicx}


% 标题与作者信息（可自行修改）
\title{Solution for Chapter 4}
\author{Zimo Guo}
\date{}

\begin{document}

\maketitle

% 使用 \noindent 去掉段落缩进，让题号更清晰
\noindent \textbf{4.1}

\begin{align*}
    E_{S \sim D^{m}}[\hat{R}_{S}(h_{S}^{ERM})] &= E_{S \sim D^{m}}[\inf \limits_{h \in \mathcal{H}} \hat{R}_{S}(h)] \\
    &\le E_{S \sim D^{m}}[\hat{R}_{S}(h)], h \in \mathcal{H} \\
    &= R(h), h \in \mathcal{H} \\
    &\le \inf \limits_{h \in \mathcal{H}}R(h)
\end{align*}

\begin{align*}
    E_{S \sim D^{m}}[R(h_{S}^{ERM})]
    &\ge E_{S \sim D^{m}}[\inf\limits_{h \in \mathcal{H}}R(h)] \\
    &= \inf\limits_{h \in \mathcal{H}}R(h) 
\end{align*}

\vspace{1em}  % 题目之间的垂直间距

\noindent \textbf{4.2}

\begin{align*}
    \Phi(u_{1}) + \Phi(u_{2}) &= (1+u_{1})^2 + (1+u_{2})^2 \\
    &= u_{1}^2 + u_{2}^2 + 2(u_{1} + u_{2}) + 2 \\
    &\ge \frac{1}{2} (u_{1} + u_{2})^2 + 2(u_{1} + u_{2}) + 2 \\
    &= 2(1 + \frac{u_{1}+u_{2}}{2})^2 \\
    &= 2\Phi(\frac{u_{1}+u_{2}}{2})
\end{align*}

So, $\Phi(u)$ is convex and non-decreasing when $u \ge -1$. Put $s=2$ and $c=\frac{1}{2}$ into
theorem 4.7 and get the objective inequality.

\vspace{1em}  % 题目之间的垂直间距

\noindent \textbf{4.3}

$\Phi(u)$ is convex and non-decreasing when $u \ge -1$. Put $s=2$ and $c=\frac{1}{2}$ into
theorem 4.7 and get the objective inequality.

\vspace{1em}  % 题目之间的垂直间距

\noindent \textbf{4.4}

\begin{enumerate}[label=(\alph*), left=0pt]
    \item 
    $\eta(x) = P[y=+1 | x]$

    Bayes classifier is a function that get $+1$ when $\eta(x) \ge \frac{1}{2}$, $-1$ when $\eta(x) <
    \frac{1}{2}$. 

    Bayes scoring function $h^{*}(x) = \eta(x) - \frac{1}{2}$.

    \item 
    \begin{align*}
        R(h) - R^{*} &= 2 E_{x \sim D_{x}}[|h^{*}(x)| \mathbb{1}_{h(x)h^{*}(x) \le 0}]
    \end{align*}
\end{enumerate}

\end{document}